\documentclass[11pt,a4paper]{article}

% -------------------- Packages --------------------
\usepackage[
top=2cm,
bottom=2cm,
left=2cm,
right=2cm
]{geometry}

\usepackage{titlesec}
\usepackage{enumitem}
\usepackage{hyperref}
\usepackage{setspace}
\usepackage{graphicx}
\usepackage{amsmath}
\usepackage{amssymb}

% -------------------- Persian --------------------
\usepackage{xepersian}

\settextfont{XB Niloofar}
\setlatintextfont{TeX Gyre Termes}

\setstretch{1.1}

% -------------------- Lists --------------------
\setlist[itemize]{
	noitemsep,
	topsep=2pt
}

\setlist[enumerate]{
	noitemsep,
	topsep=2pt
}

% -------------------- Section formatting --------------------
\titleformat{\section}
{\large\bfseries}
{\thesection}
{0.7em}
{}

\titleformat{\subsection}
{\normalsize\bfseries}
{\thesubsection}
{0.7em}
{}

\titleformat{\subsubsection}
{\normalsize\bfseries}
{\thesubsubsection}
{0.7em}
{}

% -------------------- Hyperref --------------------
\hypersetup{
	colorlinks=true,
	linkcolor=blue,
	urlcolor=blue,
	citecolor=blue
}

% ==================================================
%                   DOCUMENT
% ==================================================

\begin{document}
	
	\begin{center}
		{\LARGE \textbf{پردازش تصویر دیجیتال}}\\[0.4cm]
		{\Large \textbf{جلسه اول}}\\[0.3cm]
		{\large از «نور» تا «تصویر دیجیتال»}\\[0.5cm]
		\textit{وقتی یک تصویر در نهایت فقط جدولی از اعداد است، «پردازش تصویر» یعنی چه؟}
	\end{center}
	
	\vspace{0.3cm}
	
	% ==================================================
	\section{اهداف جلسه}
	
	در پایان این جلسه انتظار می‌رود دانشجو بتواند:
	
	\begin{itemize}
		\item \textbf{تصویر دیجیتال} را با زبان دقیق ریاضی تعریف کند.
		\item تفاوت \textbf{پردازش تصویر، تحلیل تصویر و بینایی ماشین} را توضیح دهد.
		\item نقش \textbf{طیف الکترومغناطیسی} را در منابع تصویر بیان کند و چند کاربرد برای هر ناحیه‌ی آن نام ببرد.
		\item \textbf{گام‌های اساسی} یک سامانه‌ی پردازش تصویر را نام ببرد و نقش هر گام را توضیح دهد.
		\item ساختار کلی \textbf{بینایی انسان} و تفاوت «شدت واقعی» با «روشنایی ادراک‌شده» را درک کند.
		\item \textbf{مدل روشنایی--بازتاب} را برای تشکیل تصویر بنویسد و تفسیر کند.
		\item بین \textbf{نمونه‌برداری} و \textbf{کوانتیزاسیون} تمایز قائل شود.
		\item حجم ذخیره‌سازی یک تصویر یا ویدئو را محاسبه کند.
		\item اثر \textbf{وضوح فضایی} و \textbf{وضوح شدت} را بر کیفیت تصویر توضیح دهد.
		\item یک عملیات ساده‌ی نقطه‌ای مانند \textbf{منفی تصویر} را روی ماتریس تصویر اجرا کند.
	\end{itemize}
	
	\vspace{0.3cm}
	
	\noindent
	\textbf{ایده‌ی مرکزی جلسه:}
	
	\begin{quote}
		هر تصویر دیجیتال در نهایت مجموعه‌ای از اعداد است. پردازش تصویر با این سؤال شروع می‌شود:
		
		\begin{center}
			\textbf{«با این اعداد چه کنیم تا اطلاعات مفید آشکار شود؟»}
		\end{center}
		
		دشواری اصلی از آنجا آغاز می‌شود که حجم داده‌ها زیاد است، اطلاعات در فرایند تصویربرداری
		از دست می‌رود، نویز وجود دارد و «تصویر خوب» همیشه تعریف یکتایی ندارد.
	\end{quote}
	
	
	% ==================================================
	\section{شروع بحث: تصویر در زندگی و مهندسی}
	
	بینایی پیشرفته‌ترین حس انسان است و بخش بزرگی از اطلاعاتی که از جهان می‌گیریم از راه
	تصویر است. به همین دلیل پردازش تصویر در بسیاری از حوزه‌ها دیده می‌شود.
	
	وقتی می‌پرسیم:
	
	\begin{itemize}
		\item چگونه شکستگی استخوان را در یک تصویر پرتو ایکس آشکارتر کنیم؟
		\item چگونه از روی تصویر ماهواره‌ای میزان پوشش گیاهی یا گسترش آتش‌سوزی را برآورد کنیم؟
		\item چگونه قطعه‌ی معیوب را روی خط تولید به‌طور خودکار تشخیص دهیم؟
		\item چگونه پلاک یک خودرو را از تصویر بخوانیم؟
		\item چگونه یک تصویر نویزی یا تار را بهبود دهیم؟
		\item چگونه یک تصویر یا ویدئو را با کمترین حجم ذخیره یا ارسال کنیم؟
		\item چگونه یک خودروی خودران بفهمد جلوی آن عابر پیاده است؟
	\end{itemize}
	
	در واقع با مسئله‌ای از پردازش تصویر روبه‌رو هستیم.
	
	یک نکته‌ی مهم این است که انسان فقط ناحیه‌ی بسیار باریکی از طیف الکترومغناطیسی
	(نور مرئی) را می‌بیند، اما \textbf{ماشین‌ها} می‌توانند تصویری از تقریباً همه‌ی طیف،
	از پرتوهای گاما تا امواج رادیویی، بسازند. افزون بر این، ماشین می‌تواند با تصویرهایی کار کند که
	اساساً با چشم انسان مرتبط نیستند؛ مانند تصویر فراصوت، تصویر میکروسکوپ الکترونی
	و تصویرهای تولیدشده با کامپیوتر.
	
	
	% ==================================================
	\section{تصویر چیست؟}
	
	از نظر ریاضی، یک تصویر را می‌توان تابعی دو متغیره در نظر گرفت:
	
	\[
	f(x,y)
	\]
	
	که در آن $x$ و $y$ مختصات فضایی هستند و مقدار $f$ در هر نقطه را
	\textbf{شدت} یا \lr{Intensity} (در تصویر سطح خاکستری: \textbf{سطح خاکستری}) آن نقطه می‌نامیم.
	
	اگر $x$ و $y$ و همچنین مقدار $f$ اعدادی \textbf{متناهی و گسسته} باشند، تصویر را
	\textbf{دیجیتال} می‌نامیم.
	
	یک تصویر دیجیتال با $M$ سطر و $N$ ستون را می‌توان به شکل ماتریس زیر نوشت:
	
	\[
	f(x,y)=
	\begin{bmatrix}
		f(0,0) & f(0,1) & \cdots & f(0,N-1)\\
		f(1,0) & f(1,1) & \cdots & f(1,N-1)\\
		\vdots & \vdots & \ddots & \vdots\\
		f(M-1,0) & f(M-1,1) & \cdots & f(M-1,N-1)
	\end{bmatrix}.
	\]
	
	هر عنصر این ماتریس یک \textbf{پیکسل} (\lr{Pixel} یا \lr{Picture Element}) نام دارد.
	
	چند نکته‌ی قراردادی:
	
	\begin{itemize}
		\item مبدأ مختصات معمولاً در \textbf{گوشه‌ی بالا و چپ} تصویر است.
		\item $x$ شماره‌ی سطر و $y$ شماره‌ی ستون است؛ پس $x$ به سمت پایین و $y$ به سمت راست افزایش می‌یابد.
		\item در تصویر رنگی، به‌جای یک عدد برای هر پیکسل، چند عدد (مثلاً سه مؤلفه‌ی \lr{R}، \lr{G} و \lr{B}) داریم.
	\end{itemize}
	
	\begin{quote}
		\textbf{از دید الگوریتمی، تصویر دیجیتال یک آرایه‌ی عددی است؛
			معنا و مفهوم تصویر را ما و الگوریتم‌هایمان به آن می‌دهیم.}
	\end{quote}
	
	
	% ==================================================
	\section{پردازش تصویر، تحلیل تصویر و بینایی ماشین}
	
	این سه اصطلاح اغلب با هم اشتباه گرفته می‌شوند. یک تقسیم‌بندی مفید، سه سطح پردازش است:
	
	\begin{enumerate}
		\item \textbf{پردازش سطح پایین:} ورودی و خروجی هر دو \textbf{تصویر} هستند.
		مثال: کاهش نویز، بهبود کنتراست، تیز کردن تصویر.
		\item \textbf{پردازش سطح میانی:} ورودی تصویر است، اما خروجی \textbf{ویژگی‌ها یا اجزایی استخراج‌شده از تصویر} است.
		مثال: لبه‌ها، کانتورها، نواحی (قطعه‌بندی) و اندازه‌گیری آنها.
		\item \textbf{پردازش سطح بالا:} «معنا بخشیدن» به مجموعه‌ای از اشیای تشخیص‌داده‌شده
		و انجام کارکردهای مربوط به ادراک.
		مثال: تشخیص هویت از روی چهره، فهمیدن صحنه‌ی ترافیکی.
	\end{enumerate}
	
	مرز این سه سطح کاملاً دقیق نیست، ولی برای فهم ساختار درس بسیار مفید است.
	
	مثال: \textbf{خواندن خودکار متن} (\lr{OCR}) را در نظر بگیرید:
	
	\[
	\boxed{
		\text{تصویر صفحه}
		\xrightarrow{\;\text{کاهش نویز}\;}
		\text{تصویر بهتر}
		\xrightarrow{\;\text{جداسازی حروف}\;}
		\text{نواحی حروف}
		\xrightarrow{\;\text{تشخیص}\;}
		\text{متن}
	}
	\]
	
	مرحله‌ی اول پردازش سطح پایین، مرحله‌ی دوم سطح میانی و مرحله‌ی سوم سطح بالاست.
	
	در این درس بیشتر روی سطح پایین و سطح میانی تمرکز می‌کنیم؛
	سطح بالا و مباحث وابسته به یادگیری ماشین، موضوع درس‌های دیگرند.
	
	
	% ==================================================
	\section{نگاهی کوتاه به تاریخچه}
	
	\begin{itemize}
		\item \textbf{دهه‌ی ۱۹۲۰ میلادی:} یکی از نخستین کاربردهای تصویر دیجیتال در صنعت روزنامه بود.
		تصویرها از طریق کابل زیردریایی میان لندن و نیویورک (سامانه‌ی \lr{Bartlane}) منتقل می‌شدند
		و زمان انتقال از بیش از یک هفته به کمتر از سه ساعت رسید.
		\item \textbf{دهه‌ی ۱۹۶۰:} با پیدایش کامپیوترهای قدرتمند و برنامه‌ی فضایی،
		پردازش تصویر دیجیتال جدی شد. تصویرهای ماه که توسط کاوشگر \lr{Ranger 7} (سال ۱۹۶۴)
		فرستاده شد، در آزمایشگاه \lr{JPL} با کامپیوتر اصلاح شدند.
		\item \textbf{دهه‌ی ۱۹۷۰:} اختراع \textbf{توموگرافی کامپیوتری} (\lr{CT})
		توسط \lr{Hounsfield} و \lr{Cormack}؛ آنها در سال ۱۹۷۹ جایزه‌ی نوبل پزشکی گرفتند.
		\item \textbf{از آن زمان تاکنون:} رشد سخت‌افزار، حسگرهای ارزان، دوربین‌های دیجیتال و گوشی‌های هوشمند،
		و در سال‌های اخیر یادگیری عمیق، پردازش تصویر را به فناوری روزمره تبدیل کرده است.
	\end{itemize}
	
	
	% ==================================================
	\section{منابع تصویر و طیف الکترومغناطیسی}
	
	بیشتر تصویرهایی که با آنها کار می‌کنیم از برهم‌کنش یک نوع انرژی با ماده ساخته می‌شوند.
	مهم‌ترین نوع این انرژی، تابش الکترومغناطیسی است.
	
	انرژی هر فوتون با فرکانس موج متناسب است:
	
	\[
	E=h\nu,
	\qquad
	\lambda=\frac{c}{\nu}
	\]
	
	که در آن $h$ ثابت پلانک، $\nu$ فرکانس، $\lambda$ طول موج و $c$ سرعت نور است.
	هرچه طول موج کوتاه‌تر باشد، انرژی هر فوتون بیشتر است.
	
	نواحی اصلی طیف و نمونه‌ای از کاربرد تصویربرداری آنها:
	
	\begin{itemize}
		\item \textbf{پرتو گاما:} پزشکی هسته‌ای (مثلاً \lr{PET})، نجوم، تصویربرداری از ناحیه‌های تابشی.
		\item \textbf{پرتو ایکس:} رادیوگرافی، \lr{CT}، بازرسی صنعتی و امنیتی.
		\item \textbf{فرابنفش:} میکروسکوپی فلورسانس، بازرسی صنعتی، نجوم.
		\item \textbf{نور مرئی} (تقریباً ۴۰۰ تا ۷۰۰ نانومتر): عکاسی، میکروسکوپی، بازرسی، تشخیص چهره.
		\item \textbf{فروسرخ:} تصویربرداری حرارتی، دید در شب، سنجش از دور.
		\item \textbf{ریزموج:} رادار؛ توانایی تصویربرداری از میان ابر و در شب.
		\item \textbf{امواج رادیویی:} تصویربرداری تشدید مغناطیسی (\lr{MRI}) و نجوم رادیویی.
	\end{itemize}
	
	افزون بر انرژی الکترومغناطیسی، منابع دیگری نیز وجود دارد:
	
	\begin{itemize}
		\item \textbf{صوت:} فراصوت در پزشکی، اکتشاف زمین‌شناسی.
		\item \textbf{باریکه‌ی الکترونی:} میکروسکوپ الکترونی.
		\item \textbf{تصویرهای ساختگی:} تصویرهای گرافیک کامپیوتری و تصویرهای مصنوعی برای آزمایش الگوریتم‌ها.
	\end{itemize}
	
	نکته‌ی مهم این است که الگوریتم‌های پردازش، معمولاً \textbf{فارغ از نوع منبع} روی همان
	ماتریس عددی کار می‌کنند؛ اما تفسیر نتیجه و انتخاب روش مناسب به منبع تصویر بستگی دارد.
	
	
	% ==================================================
	\section{گام‌های اساسی در پردازش تصویر دیجیتال}
	
	یک سامانه‌ی کامل پردازش تصویر را می‌توان به گام‌های زیر تقسیم کرد:
	
	\begin{enumerate}
		\item \textbf{تصویربرداری} (\lr{Image Acquisition}): گرفتن تصویر و تبدیل آن به فرم دیجیتال.
		\item \textbf{بهبود تصویر} (\lr{Enhancement}): آشکار کردن جزئیات مورد علاقه
		یا بهتر کردن تصویر برای یک کاربرد مشخص؛ فرایندی \textbf{ذهنی و وابسته به کاربرد}.
		\item \textbf{بازسازی تصویر} (\lr{Restoration}): برگرداندن تصویر به حالت اصلی با کمک مدل ریاضی
		تخریب (مثل تاری حرکت یا نویز)؛ فرایندی \textbf{عینی‌تر} از بهبود.
		\item \textbf{پردازش تصویر رنگی} (\lr{Color Image Processing}).
		\item \textbf{موجک‌ها و پردازش چندمقیاسی} (\lr{Wavelets and Multiresolution}).
		\item \textbf{فشرده‌سازی} (\lr{Compression}): کاهش حجم ذخیره‌سازی یا پهنای باند لازم.
		\item \textbf{پردازش ریخت‌شناسی} (\lr{Morphological Processing}): استخراج مؤلفه‌های مفید
		برای نمایش شکل.
		\item \textbf{قطعه‌بندی} (\lr{Segmentation}): تقسیم تصویر به نواحی یا اشیا.
		\item \textbf{نمایش و توصیف} (\lr{Representation and Description}): تبدیل نتیجه‌ی قطعه‌بندی
		به شکلی مناسب برای پردازش کامپیوتری (مرز یا ناحیه) و استخراج ویژگی.
		\item \textbf{تشخیص اشیا} (\lr{Object Recognition}): نسبت دادن برچسب به شیء بر اساس توصیف آن.
	\end{enumerate}
	
	در پس همه‌ی این گام‌ها یک \textbf{پایگاه دانش} (\lr{Knowledge Base}) قرار دارد
	که اطلاعات ما درباره‌ی مسئله را نگه می‌دارد؛ مثلاً ناحیه‌هایی که انتظار داریم اطلاعات مهم در آن باشد.
	
	توجه کنید که خروجی گام‌های اول تا هفتم معمولاً \textbf{تصویر} است، در حالی که از گام هشتم
	به بعد خروجی \textbf{ویژگی، ساختار یا تصمیم} است.
	
	همچنین لازم نیست در هر کاربرد از همه‌ی گام‌ها استفاده شود.
	
	
	% ==================================================
	\section{اجزای یک سامانه‌ی پردازش تصویر}
	
	یک سامانه‌ی معمول پردازش تصویر دیجیتال شامل اجزای زیر است:
	
	\begin{itemize}
		\item \textbf{حسگر تصویر:} یک وسیله‌ی حساس به انرژی که سیگنال الکتریکی متناسب با انرژی دریافتی تولید می‌کند.
		\item \textbf{دیجیتایزر:} تبدیل خروجی حسگر به داده‌ی دیجیتال.
		\item \textbf{سخت‌افزار تخصصی} (مانند \lr{GPU}) برای عملیات سریع و سنگین.
		\item \textbf{کامپیوتر:} از رایانه‌ی شخصی تا ابررایانه.
		\item \textbf{نرم‌افزار:} ماژول‌ها و کتابخانه‌های پردازش تصویر.
		\item \textbf{ذخیره‌سازی:} حافظه‌ی کوتاه‌مدت، حافظه‌ی برخط و آرشیو.
		\item \textbf{نمایشگر و چاپگر.}
		\item \textbf{شبکه:} برای ارسال و دریافت تصویر؛ به‌ویژه وقتی پهنای باند محدود است، فشرده‌سازی اهمیت پیدا می‌کند.
	\end{itemize}
	
	
	% ==================================================
	\section{بینایی انسان: چرا لازم است درباره‌ی آن بدانیم؟}
	
	بسیاری از کاربردهای پردازش تصویر در نهایت به یک «بیننده‌ی انسانی» ختم می‌شوند؛
	پس باید بدانیم چشم انسان چگونه کار می‌کند و چه محدودیت‌هایی دارد.
	
	\subsection{ساختار چشم و گیرنده‌ها}
	
	روی شبکیه‌ی چشم دو نوع گیرنده‌ی نور وجود دارد:
	
	\begin{itemize}
		\item \textbf{مخروط‌ها} (\lr{Cones}): تعداد آنها در هر چشم حدود ۶ تا ۷ میلیون است،
		عمدتاً در ناحیه‌ی مرکزی شبکیه (\lr{Fovea}) قرار دارند، به رنگ حساس‌اند
		و دید با جزئیات زیاد در روشنایی کافی را ممکن می‌کنند.
		\item \textbf{میله‌ها} (\lr{Rods}): تعداد آنها حدود ۷۵ تا ۱۵۰ میلیون است،
		در سطح وسیعی از شبکیه پخش شده‌اند، به نور کم حساس‌اند و دید کلی و در نور کم را فراهم می‌کنند.
	\end{itemize}
	
	در ناحیه‌ای از شبکیه که عصب بینایی خارج می‌شود (\textbf{نقطه‌ی کور}) گیرنده وجود ندارد.
	
	\subsection{روشنایی ادراک‌شده با شدت واقعی یکی نیست}
	
	چند واقعیت مهم:
	
	\begin{itemize}
		\item سامانه‌ی بینایی انسان می‌تواند بازه‌ی بسیار بزرگی از شدت نور (از مرتبه‌ی $10^{10}$) را پوشش دهد؛
		اما در یک لحظه فقط بازه‌ی کوچکی از آن را می‌بیند (\textbf{تطبیق با روشنایی}).
		\item روشنایی ادراک‌شده تابعی \textbf{غیرخطی} از شدت واقعی است.
		کوچک‌ترین تغییر قابل تشخیص در شدت، به شدت زمینه وابسته است
		(نسبت وبر: $\Delta I_c/I$ حدوداً ثابت).
		\item \textbf{اثر ماخ} (\lr{Mach Bands}): در مرز بین دو ناحیه‌ی روشن و تاریک،
		نوارهایی روشن‌تر یا تاریک‌تر از مقدار واقعی دیده می‌شود.
		\item \textbf{تضاد هم‌زمان} (\lr{Simultaneous Contrast}): یک پچ خاکستری با شدت ثابت
		روی زمینه‌ی روشن تیره‌تر از زمینه‌ی تاریک به نظر می‌رسد.
	\end{itemize}
	
	\begin{quote}
		\textbf{مقدار یک پیکسل با «آنچه می‌بینیم» یکی نیست؛
			ادراک انسان یک اندازه‌گیری دقیق نیست.}
	\end{quote}
	
	به همین دلیل بهبود تصویر برای انسان نمی‌تواند فقط به عدد و فرمول تکیه کند
	و در بسیاری از موارد، ارزیابی نهایی ذهنی است.
	
	
	% ==================================================
	\section{مدل تشکیل تصویر: روشنایی و بازتاب}
	
	تصویری که حسگر دریافت می‌کند، ناشی از دو عامل است:
	
	\begin{itemize}
		\item \textbf{روشنایی} (\lr{Illumination}) $i(x,y)$: مقدار نوری که به صحنه می‌تابد.
		\item \textbf{بازتاب} (\lr{Reflectance}) $r(x,y)$: بخشی از نور که شیء بازمی‌تاباند.
	\end{itemize}
	
	مدل ساده‌ی تشکیل تصویر:
	
	\[
	f(x,y)=i(x,y)\,r(x,y)
	\]
	
	با شرایط:
	
	\[
	0<i(x,y)<\infty,
	\qquad
	0<r(x,y)<1.
	\]
	
	مقدار $r=0$ یعنی جذب کامل و $r=1$ یعنی بازتاب کامل.
	
	برای درک بزرگی این اعداد:
	
	\begin{itemize}
		\item روشنایی یک روز آفتابی در حدود $90000$ و محیط اداری در حدود $1000$
		و شب مهتابی در حدود $0.1$ لومن بر متر مربع است.
		\item بازتاب مخمل سیاه حدود $0.01$ و برف حدود $0.93$ است.
	\end{itemize}
	
	نتیجه‌ی مهم این است که یک شیء با بازتاب ثابت، بسته به روشنایی صحنه، می‌تواند مقدارهای بسیار
	متفاوتی در تصویر ایجاد کند. این یکی از دلایل دشواری بسیاری از مسائل پردازش تصویر است.
	
	مثال: اگر $i=1000$ و $r=0.8$ باشد:
	
	\[
	f=1000\times0.8=800.
	\]
	
	و اگر روشنایی نصف شود، مقدار تصویر هم نصف خواهد شد، در حالی که خود شیء تغییری نکرده است.
	
	
	% ==================================================
	\section{تصویربرداری: از نور تا ولتاژ}
	
	حسگر تصویر انرژی ورودی را به یک سیگنال الکتریکی (مثلاً ولتاژ) تبدیل می‌کند.
	سه آرایش رایج داریم:
	
	\begin{enumerate}
		\item \textbf{حسگر منفرد:} برای ساخت یک تصویر دوبعدی باید حسگر یا نمونه به‌صورت مکانیکی حرکت کند.
		ارزان است اما کند.
		\item \textbf{نوار حسگر:} یک ردیف حسگر؛ حرکت در جهت دیگر تصویر را کامل می‌کند
		(مانند اسکنر و بسیاری از دستگاه‌های سنجش از دور).
		\item \textbf{آرایه‌ی حسگر:} شبکه‌ی دوبعدی از حسگرها؛ شکل رایج در دوربین‌های دیجیتال
		(از نوع \lr{CCD} یا \lr{CMOS}).
	\end{enumerate}
	
	هر عنصر حسگر، نوری را که در یک ناحیه‌ی کوچک دریافت می‌کند، به یک مقدار پیوسته تبدیل می‌کند.
	این خروجی هنوز \textbf{پیوسته} است و برای ورود به کامپیوتر باید دیجیتال شود.
	
	
	% ==================================================
	\section{نمونه‌برداری و کوانتیزاسیون}
	
	برای تبدیل یک تصویر پیوسته $f(x,y)$ به تصویر دیجیتال باید دو کار انجام شود:
	
	\begin{enumerate}
		\item \textbf{نمونه‌برداری} (\lr{Sampling}): گسسته کردن \textbf{مختصات}
		$x$ و $y$؛ یعنی تصمیم می‌گیریم تصویر را در چند نقطه اندازه بگیریم.
		\item \textbf{کوانتیزاسیون} (\lr{Quantization}): گسسته کردن \textbf{مقدار شدت}؛
		یعنی تصمیم می‌گیریم مقدار هر نمونه را با چند سطح مجاز نمایش دهیم.
	\end{enumerate}
	
	\[
	\boxed{
		\text{نمونه‌برداری}\;\Rightarrow\;\text{چند پیکسل؟}\;(M\times N)
		\qquad
		\text{کوانتیزاسیون}\;\Rightarrow\;\text{چند سطح شدت؟}\;(L)
	}
	\]
	
	معمولاً تعداد سطوح شدت را توانی از دو انتخاب می‌کنیم:
	
	\[
	L=2^{k}.
	\]
	
	که در آن $k$ تعداد بیت لازم برای هر پیکسل است. برای مثال:
	
	\[
	k=1\Rightarrow L=2,
	\qquad
	k=4\Rightarrow L=16,
	\qquad
	k=8\Rightarrow L=256.
	\]
	
	در تصویر $8$ بیتی، شدت‌ها اعداد صحیح بازه‌ی $[0,255]$ هستند؛
	معمولاً $0$ سیاه و $255$ سفید است.
	
	چند نکته:
	
	\begin{itemize}
		\item نمونه‌برداری درباره‌ی \textbf{مکان} است و کوانتیزاسیون درباره‌ی \textbf{مقدار}.
		\item هر دو باعث از دست رفتن اطلاعات می‌شوند؛ تبدیل برگشت‌ناپذیر است.
		\item انتخاب $M$ و $N$ و $L$ یک \textbf{مصالحه} بین کیفیت و هزینه است.
		\item این‌که چقدر نمونه‌برداری کافی است، به \textbf{قضیه‌ی نمونه‌برداری} و پدیده‌ی
		\lr{Aliasing} مربوط می‌شود که در جلسات آینده بررسی می‌شود.
	\end{itemize}
	
	
	% ==================================================
	\section{حجم ذخیره‌سازی تصویر}
	
	تعداد بیت لازم برای ذخیره‌ی یک تصویر دیجیتال $M\times N$ با $k$ بیت در هر پیکسل برابر است با:
	
	\[
	b=M\times N\times k.
	\]
	
	\textbf{مثال ۱:} تصویر خاکستری $1024\times1024$ با $8$ بیت برای هر پیکسل:
	
	\[
	b=2^{10}\times2^{10}\times8=2^{23}\ \text{بیت}=2^{20}\ \text{بایت}\approx 1\ \text{مگابایت}.
	\]
	
	\textbf{مثال ۲:} یک تصویر رنگی $3000\times2000$ با سه کانال و $8$ بیت برای هر کانال:
	
	\[
	3000\times2000\times3=18{,}000{,}000\ \text{بایت}\approx 18\ \text{مگابایت}.
	\]
	
	\textbf{مثال ۳:} یک ویدئوی رنگی $1920\times1080$ با $30$ فریم در ثانیه، بدون فشرده‌سازی:
	
	\[
	1920\times1080\times3=6{,}220{,}800\ \text{بایت در هر فریم}
	\]
	
	\[
	6{,}220{,}800\times30=186{,}624{,}000\ \text{بایت در ثانیه}
	\approx 1.5\ \text{گیگابیت بر ثانیه}.
	\]
	
	این عدد توضیح می‌دهد که چرا \textbf{فشرده‌سازی} یکی از گام‌های اساسی پردازش تصویر است
	و چرا حجم داده، خودش یک چالش مهندسی محسوب می‌شود.
	
	
	% ==================================================
	\section{وضوح فضایی و وضوح شدت}
	
	\subsection{وضوح فضایی}
	
	\textbf{وضوح فضایی} (\lr{Spatial Resolution}) کوچک‌ترین جزئیات قابل تشخیص در تصویر است.
	معمولاً با «تعداد نقطه در واحد طول» مانند \lr{dpi} (نقطه در اینچ) بیان می‌شود.
	
	اگر تعداد پیکسل‌ها را کم کنیم (مثلاً با نگه داشتن یک پیکسل از هر چهار پیکسل)،
	اندازه‌ی فیزیکی تصویر ثابت است، اما جزئیات از بین می‌روند و در حالت شدید،
	ساختاری شبیه صفحه‌ی شطرنجی (\lr{Checkerboard Effect}) دیده می‌شود.
	
	\subsection{وضوح شدت}
	
	\textbf{وضوح شدت} (\lr{Intensity Resolution}) کوچک‌ترین تغییر قابل تشخیص در شدت است،
	که عملاً به تعداد سطوح شدت، یعنی $L=2^k$، وابسته است.
	
	اگر تعداد سطوح را کم کنیم، در نواحی هموار و تدریجی (مانند آسمان یا پوست) خطوط ریز و
	مصنوعی ظاهر می‌شوند که \textbf{کانتور کاذب} (\lr{False Contouring}) نام دارد.
	
	\subsection{کدام مهم‌تر است؟}
	
	جواب ثابتی وجود ندارد.
	
	\begin{itemize}
		\item در تصویری با جزئیات زیاد (مثل یک صحنه‌ی شلوغ)، کاهش پیکسل‌ها بیشتر آسیب می‌زند.
		\item در تصویری با نواحی هموار، کاهش سطوح شدت بیشتر آسیب می‌زند.
	\end{itemize}
	
	\begin{quote}
		\textbf{«تصویر بزرگ‌تر» یا «بیت بیشتر» همیشه به معنی «تصویر بهتر» نیست؛
			باید دید کاربرد چه می‌خواهد.}
	\end{quote}
	
	
	% ==================================================
	\section{عملیات روی تصویر: نخستین نگاه}
	
	هر پردازش روی تصویر را می‌توان به شکل زیر نشان داد:
	
	\[
	g(x,y)=T\big[f(x,y)\big]
	\]
	
	که در آن $f$ تصویر ورودی، $g$ تصویر خروجی و $T$ یک عملگر است.
	
	ساده‌ترین حالت وقتی است که $T$ فقط به مقدار \textbf{همان پیکسل} نگاه کند
	(\textbf{عملیات نقطه‌ای}). اگر $r$ شدت ورودی و $s$ شدت خروجی باشد:
	
	\[
	s=T(r).
	\]
	
	یک مثال بسیار ساده، \textbf{منفی تصویر} است:
	
	\[
	s=L-1-r.
	\]
	
	این عملیات در تصویر پزشکی گاهی برای دیدن جزئیات سفید روی زمینه‌ی تیره به کار می‌رود.
	
	نمونه‌ی دیگر عملیات، \textbf{عملیات همسایگی} است که برای تعیین مقدار خروجی در یک پیکسل،
	از مقدار پیکسل‌های اطراف آن هم استفاده می‌کند (مثلاً میانگین‌گیری برای کاهش نویز).
	این موضوع را در جلسات بعد با جزئیات بررسی می‌کنیم.
	
	
	% ==================================================
	\section{یک مثال کامل: منفی یک تصویر $3$ بیتی}
	
	تصویر $4\times4$ زیر را با $k=3$ در نظر بگیرید؛ پس $L=8$ و شدت‌ها در بازه‌ی $[0,7]$ هستند:
	
	\[
	f=
	\begin{bmatrix}
		2 & 2 & 5 & 7\\
		1 & 3 & 6 & 7\\
		0 & 3 & 4 & 5\\
		0 & 1 & 2 & 4
	\end{bmatrix}.
	\]
	
	با رابطه‌ی $s=L-1-r=7-r$:
	
	\[
	g=
	\begin{bmatrix}
		5 & 5 & 2 & 0\\
		6 & 4 & 1 & 0\\
		7 & 4 & 3 & 2\\
		7 & 6 & 5 & 3
	\end{bmatrix}.
	\]
	
	چند مشاهده:
	
	\begin{itemize}
		\item حجم این تصویر برابر $4\times4\times3=48$ بیت است.
		\item پیکسل‌های روشن به تاریک و پیکسل‌های تاریک به روشن تبدیل شدند.
		\item اگر همین عملیات را یک بار دیگر اجرا کنیم، به تصویر اصلی برمی‌گردیم؛ زیرا
		$7-(7-r)=r$.
	\end{itemize}
	
	
	% ==================================================
	\section{چرا پردازش تصویر دشوار است؟}
	
	در جلسه‌های اولیه ممکن است به نظر برسد که تصویر فقط یک آرایه‌ی عددی است و کار با آن ساده است.
	اما دشواری‌ها واقعی‌اند:
	
	\begin{itemize}
		\item \textbf{حجم زیاد داده:} همان‌طور که دیدیم، چند ثانیه ویدئو می‌تواند چندین گیگابایت داده باشد.
		\item \textbf{از دست رفتن اطلاعات:} تصویربرداری، سه‌بعدی را به دوبعدی تبدیل می‌کند
		و نمونه‌برداری و کوانتیزاسیون هم اطلاعات را کم می‌کنند.
		\item \textbf{نویز:} هر حسگر واقعی مقداری خطا دارد.
		\item \textbf{تغییر شرایط روشنایی:} طبق مدل $f=i\,r$، یک شیء ثابت می‌تواند تصویرهای کاملاً متفاوتی بسازد.
		\item \textbf{ابهام:} چند صحنه‌ی متفاوت می‌توانند تصویر مشابهی بسازند.
		\item \textbf{نبود معیار یکتا برای «بهتر بودن»:} یک تصویر ممکن است برای انسان خوب و برای الگوریتم بد باشد یا برعکس.
	\end{itemize}
	
	بنابراین پیش از انتخاب هر روش، باید دقیق مشخص کنیم که
	\textbf{هدف نهایی چیست} و خروجی چه کسی یا چه الگوریتمی را قرار است راضی کند.
	
	
	% ==================================================
	\section{جمع‌بندی مفهومی جلسه}
	
	تا اینجا مسیر زیر را طی کردیم:
	
	\[
	\boxed{
		\text{صحنه}
		\rightarrow
		\text{روشنایی و بازتاب}
		\rightarrow
		\text{حسگر}
		\rightarrow
		\text{نمونه‌برداری و کوانتیزاسیون}
		\rightarrow
		\text{ماتریس عددی}
		\rightarrow
		\text{پردازش}
	}
	\]
	
	و مهم‌ترین اصطلاحات جلسه عبارت‌اند از:
	
	\begin{itemize}
		\item \textbf{تصویر دیجیتال}:
		تابع $f(x,y)$ که مختصات و مقدار آن گسسته و متناهی‌اند.
		
		\item \textbf{پیکسل}:
		یک عنصر از آرایه‌ی تصویر.
		
		\item \textbf{شدت / سطح خاکستری}:
		مقدار تصویر در یک پیکسل.
		
		\item \textbf{نمونه‌برداری}:
		گسسته کردن مختصات فضایی.
		
		\item \textbf{کوانتیزاسیون}:
		گسسته کردن مقدار شدت.
		
		\item \textbf{وضوح فضایی}:
		میزان جزئیات مکانی قابل تشخیص.
		
		\item \textbf{وضوح شدت}:
		تعداد سطوح شدت قابل نمایش.
		
		\item \textbf{مدل روشنایی--بازتاب}:
		$f(x,y)=i(x,y)\,r(x,y)$.
		
		\item \textbf{پردازش سطح پایین، میانی و بالا}:
		به ترتیب تصویر به تصویر، تصویر به ویژگی، و ویژگی به معنا.
		
		\item \textbf{بهبود و بازسازی}:
		اولی ذهنی و وابسته به کاربرد، دومی بر پایه‌ی مدل تخریب.
	\end{itemize}
	
	
	% ==================================================
	\section{تمرین کلاسی ۱: سطح پردازش}
	
	برای هر کار زیر مشخص کنید که در کدام سطح قرار می‌گیرد
	(پایین، میانی یا بالا) و ورودی و خروجی آن چیست:
	
	\begin{enumerate}
		\item حذف نویز از یک عکس گرفته‌شده با گوشی؛
		\item استخراج لبه‌های یک تصویر؛
		\item شمارش سلول‌ها در یک تصویر میکروسکوپی؛
		\item خواندن پلاک خودرو از تصویر دوربین ترافیکی؛
		\item بهتر کردن کنتراست یک تصویر پرتو ایکس؛
		\item فشرده‌سازی تصویر برای ارسال از طریق شبکه؛
		\item باز کردن قفل گوشی با تشخیص چهره.
	\end{enumerate}
	
	
	% ==================================================
	\section{تمرین کلاسی ۲: منبع تصویر مناسب}
	
	برای هر کاربرد، یک ناحیه‌ی طیف (یا یک روش تصویربرداری) پیشنهاد دهید
	و دلیل انتخاب خود را در یک جمله بنویسید:
	
	\begin{enumerate}
		\item تشخیص شکستگی استخوان؛
		\item دید در شب؛
		\item پایش وضعیت جنگل‌ها از ماهواره؛
		\item تصویربرداری از جنین؛
		\item مشاهده‌ی ساختار نرم مغز؛
		\item بازرسی بدون تخریب یک قطعه‌ی فلزی؛
		\item تصویربرداری از آسمان در شب ابری.
	\end{enumerate}
	
	
	% ==================================================
	\section{تمرین ریاضی}
	
	\begin{enumerate}
		\item یک تصویر خاکستری $512\times512$ با $64$ سطح شدت چند بیت حجم دارد؟
		(مقدار $k$ را هم مشخص کنید.)
		\item یک تصویر رنگی $3840\times2160$ با سه کانال $8$ بیتی چند مگابایت است؟
		\item یک ویدئوی $1280\times720$ رنگی با سه کانال $8$ بیتی و $25$ فریم در ثانیه
		در هر ساعت چند گیگابایت داده‌ی خام تولید می‌کند؟
		\item اگر $i(x,y)=500$ و $r(x,y)=0.2$ باشد، $f(x,y)$ را حساب کنید.
		اگر روشنایی دو برابر شود چه می‌شود؟
		\item منفی تصویر $3$ بیتی زیر را بنویسید:
		\[
		f=
		\begin{bmatrix}
			0 & 1 & 2\\
			3 & 4 & 5\\
			6 & 7 & 7
		\end{bmatrix}.
		\]
		\item نشان دهید که منفی منفی یک تصویر، همان تصویر اصلی است.
	\end{enumerate}
	
	
	% ==================================================
	\section{تمرین فکری: وضوح و کیفیت}
	
	فرض کنید یک تصویر $1000\times1000$ با $8$ بیت برای هر پیکسل داریم.
	
	\begin{enumerate}
		\item اگر ابعاد تصویر را نصف کنیم (با نگه داشتن یک پیکسل از هر چهار پیکسل)،
		تعداد پیکسل‌ها و حجم چقدر می‌شود؟
		\item اگر تعداد سطوح شدت را به $2$ کاهش دهیم، چه نوع تخریبی انتظار دارید؟
		\item برای یک تصویر چهره و برای یک تصویر آسمان صاف، کاهش کدام‌یک از دو نوع وضوح
		آسیب بیشتری می‌زند؟ چرا؟
		\item آیا می‌توان تصویر اصلی را از روی تصویر کوچک‌شده «دقیقاً» بازسازی کرد؟ چرا؟
	\end{enumerate}
	
	هدف این تمرین، بیش از محاسبه‌ی عددی، درک این نکته است که کیفیت تصویر
	بین چند عامل در تعادل است.
	
	
	% ==================================================
	\section{تمرین عملی (در نرم‌افزار)}
	
	با \lr{Python} (کتابخانه‌هایی مانند \lr{NumPy} و \lr{OpenCV}) یا \lr{MATLAB} کارهای زیر را انجام دهید:
	
	\begin{enumerate}
		\item یک تصویر خاکستری را بخوانید و اندازه، نوع داده، کمینه و بیشینه‌ی مقدار آن را چاپ کنید.
		\item تصویر را نمایش دهید و مقدار چند پیکسل را با دست بررسی کنید.
		\item تعداد سطوح شدت را به $16$، $4$ و $2$ کاهش دهید و اثر آن را مشاهده کنید.
		\item تصویر را با ضریب $2$، $4$ و $8$ کوچک کنید (با نگه داشتن یک پیکسل از هر چند پیکسل)
		و دوباره با همان اندازه‌ی اولیه نمایش دهید.
		\item منفی تصویر را بسازید.
		\item در یک گزارش کوتاه بنویسید در چه مرحله‌ای، چه نوع تخریبی ظاهر شد.
	\end{enumerate}
	
	
	% ==================================================
	\section{تمرین پیش از جلسه‌ی دوم}
	
	برای جلسه‌ی بعد، یک \textbf{کاربرد واقعی پردازش تصویر} انتخاب کنید
	و آن را در حداکثر یک صفحه توضیح دهید:
	
	\begin{enumerate}
		\item هدف نهایی کاربرد چیست؟
		\item تصویر از چه منبعی و با چه حسگری گرفته می‌شود؟
		\item خروجی سامانه تصویر است، ویژگی است یا تصمیم؟
		\item چه گام‌هایی از «گام‌های اساسی پردازش تصویر» در آن لازم است؟
		\item سخت‌ترین بخش مسئله کدام است (نویز، روشنایی، حجم داده، ...)؟
		\item یک وضعیت که ممکن است سامانه در آن شکست بخورد را توضیح دهید.
	\end{enumerate}
	
	این کاربرد می‌تواند در ادامه‌ی ترم به پروژه‌ی اصلی دانشجو تبدیل شود.
	
	
	% ==================================================
	\section{پیوند جلسه‌ی اول با ادامه‌ی درس}
	
	در این جلسه دیدیم که یک تصویر دیجیتال چیست و چگونه ساخته می‌شود.
	
	اما هنوز از آن «استفاده» نکرده‌ایم. اکنون یک ماتریس عددی داریم:
	
	\[
	\boxed{
		f(x,y),\qquad 0\leq x\leq M-1,\quad 0\leq y\leq N-1
	}
	\]
	
	و در جلسه‌ی دوم می‌پرسیم:
	
	\begin{center}
		\Large
		\textbf{چگونه می‌توانیم مقدار این پیکسل‌ها را به‌طور هدفمند تغییر دهیم؟}
	\end{center}
	
	سپس مفاهیمی مانند:
	
	\[
	\text{همسایگی و مجاورت پیکسل‌ها}
	\]
	
	\[
	\text{فاصله میان پیکسل‌ها}
	\]
	
	\[
	\text{تبدیل‌های شدتی (لگاریتمی، توانی، کشش کنتراست)}
	\]
	
	و
	
	\[
	\text{هیستوگرام تصویر}
	\]
	
	وارد بحث خواهند شد.
	
	در واقع جلسه‌ی دوم پلی خواهد بود از:
	
	\[
	\boxed{
		\text{تصویر دیجیتال}
		\longrightarrow
		\text{عملیات روی پیکسل‌ها}
		\longrightarrow
		\text{بهبود تصویر}
	}
	\]
	
	
	% ==================================================
	\section{خلاصه‌ی نهایی جلسه}
	
	در این جلسه دیدیم که پردازش تصویر فقط «زیباتر کردن عکس» نیست؛ بلکه
	چارچوبی برای استخراج اطلاعات از داده‌هایی است که از اندازه‌گیری نور و انرژی به دست می‌آیند.
	
	هر تصویر دیجیتال را می‌توان با سه عنصر اصلی توصیف کرد:
	
	\[
	\boxed{
		M\times N,
		\quad
		L=2^{k},
		\quad
		f(x,y)
	}
	\]
	
	که در آن:
	
	\begin{itemize}
		\item $M\times N$ تعداد پیکسل‌ها (وضوح فضایی) است؛
		\item $L=2^k$ تعداد سطوح شدت (وضوح شدت) است؛
		\item $f(x,y)$ مقدار شدت در هر پیکسل است.
	\end{itemize}
	
	حجم داده برابر است با:
	
	\[
	b=M\times N\times k.
	\]
	
	مدل ساده‌ی تشکیل تصویر نیز چنین است:
	
	\[
	f(x,y)=i(x,y)\,r(x,y).
	\]
	
	مشکل اصلی زمانی پدیدار می‌شود که:
	
	\begin{itemize}
		\item حجم داده بسیار زیاد باشد؛
		\item اطلاعات در تصویربرداری از دست رفته باشند؛
		\item نویز و تغییر روشنایی وجود داشته باشد؛
		\item و معیار «خوب بودن» نتیجه، به کاربرد وابسته باشد.
	\end{itemize}
	
	برای طراحی یک سامانه‌ی پردازش تصویر خوب باید دائماً به سه سؤال برگردیم:
	
	\[
	\boxed{
		\begin{array}{c}
			\text{تصویر چگونه ساخته شده است؟}\\
			\text{چه اطلاعاتی در آن مانده و چه چیزی از دست رفته است؟}\\
			\text{چه چیزی را می‌خواهیم از آن به دست آوریم؟}
		\end{array}
	}
	\]
	
	سؤال اول و دوم را امروز آغاز کردیم.
	
	پرداختن به سؤال سوم، با عملیات روی پیکسل‌ها، از جلسه‌ی دوم آغاز می‌شود.
	
	
	% ==================================================
	\section{منابع}
	
	\begin{enumerate}
		
		\item
		Rafael C. Gonzalez and Richard E. Woods,
		\textit{Digital Image Processing},
		4th Edition, Pearson, 2018
		(فصل اول: \lr{Introduction}؛ فصل دوم: \lr{Digital Image Fundamentals}).
		
		\item
		Rafael C. Gonzalez, Richard E. Woods, and Steven L. Eddins,
		\textit{Digital Image Processing Using MATLAB},
		Gatesmark Publishing.
		
		\item
		Anil K. Jain,
		\textit{Fundamentals of Digital Image Processing},
		Prentice Hall, 1989.
		
		\item
		Richard Szeliski,
		\textit{Computer Vision: Algorithms and Applications},
		Springer.
		
	\end{enumerate}
	
	
	% ==================================================
	\section*{یک جمله برای به خاطر سپردن}
	
	\begin{center}
		\Large
		\textbf{
			تصویر دیجیتال یعنی نمونه‌برداری از مکان و از شدت نور؛
			پردازش تصویر یعنی تصمیم گرفتن این‌که با این اعداد چه کنیم
			تا اطلاعات مفید آشکار شود.
		}
	\end{center}
	
\end{document}
